\section{Sensor de flujo}

El sensor mide el caudal real administrado por la bomba y emite una lectura cada segundo.

\subsection{Definición formal}

\begin{align*}
  X &= \{\,(c,\;0)\;\mid\; c\in[0,200]\,\} \\[4pt]
  Y &= \{\,(c,\;0)\;\mid\; c\in[0,200]\,\} \\[4pt]
  S &= \{\,(c,\;\sigma)\;\mid\; c\in[0,200],\;\sigma\in\mathbb{R}_0^{+}\,\}
\end{align*}

El estado $s = (c,\,\sigma)$ contiene el último caudal recibido del actuador y el tiempo restante $\sigma$ hasta la próxima medición.

\paragraph{Transición externa.}
Actualiza el caudal medido sin modificar el tiempo restante:
\[
  \dext\bigl((c,\,\sigma),\;e,\;(x,0)\bigr)
  \;=\; (x,\;\sigma - e).
\]

\paragraph{Transición interna.}
Emite la medición y programa la siguiente un segundo después:
\[
  \dint(c,\,\sigma) \;=\; (c,\;1).
\]

\paragraph{Función de salida.}
\[
  \lambda(c,\,\sigma) \;=\; (c,\;0).
\]

\paragraph{Avance temporal.}
\[
  \ta(c,\,\sigma) \;=\; \sigma.
\]

\paragraph{Estado inicial.}
\[
  s_0 \;=\; (0,\;\infty).
\]
El sensor permanece pasivo hasta recibir la primera lectura del actuador.

\subsection{Diagrama de estados}

\begin{figure}[htbp]
\centering
\begin{tikzpicture}
  \node[draw, rectangle] (act) at (-5,0) {Actuador de la bomba};
  \node[draw, circle]    (sensor) at (0,0) {$(c,\,\sigma)$};
  \node[draw, rectangle] (ctrl)   at (5,0) {Controlador};

  \draw[-{Stealth}] (act) -- (sensor)
    node[midway, above] {$\lambda=(x,0)$};

  \draw[-{Stealth}] (sensor) -- (ctrl)
    node[midway, above] {$\lambda = (c,\,0)$};

  \draw[-{Stealth}]
    (sensor) edge[loop above]
    node {$\delta_{\text{int}}(c,\,\sigma) = (c,\,1)$}
    (sensor);

  \draw[-{Stealth}]
    (sensor) edge[loop below]
    node {$\delta_{\text{ext}}\bigl((c,\,\sigma),\,e,\,(x,0)\bigr) = (x,\,\sigma-e)$}
    (sensor);
\end{tikzpicture}
\caption{Diagrama de estados del Sensor de Flujo.}
\label{fig:diag_sensor}
\end{figure}